\chapter{Final Project Outcomes and Evaluation}

This chapter summarizes the outcomes achieved by the SiteX project after the completion of the HeteroGAT-based development cycle. The project delivered a working end-to-end platform that combined graph-based prediction, contextualized scoring, and a streamlined UI workflow for cafe site selection in the Kathmandu Valley.

\section*{1. Successful Training and Deployment of HeteroGAT}
A major outcome of this project was the successful training and deployment of the HeteroGAT model as the production prediction engine. The graph-based architecture was integrated into the backend and used for real-time location suitability inference, replacing the earlier baseline-only scoring approach.

\section*{2. Production Integration and Contextualized Scoring}
The completed system now combines the HeteroGAT location score with audience-fit, road accessibility, and competition-penalty factors. This contextualized scoring approach produced a more realistic and business-aware suitability recommendation for each user profile, rather than a single raw model output.

\section*{3. Sentiment-Enriched and Traffic-Aware Analysis}
The project also delivered a sentiment-enriched analysis pipeline that incorporated review quality and foot-traffic signals into the overall site evaluation. The documentation records that the sentiment pipeline successfully analyzed 209,411 reviews using VADER sentiment analysis with a 0\% error rate, demonstrating a high-confidence and scalable analysis workflow.

\section*{4. UI v2.1 Delivery and User Workflow}
The completed interface introduced a project-dashboard landing experience, a six-step Requirements Input Form, and a seven-tab secondary analysis workspace. These features enabled users to move from business-profile entry to location review in a structured and interactive manner, improving both usability and decision support.

\section*{5. Evaluation and Project Readiness}
Overall, the project achieved its core objective of delivering a deployable location intelligence platform for cafe site selection. The completed system is now positioned not only as a research prototype but also as a practical decision-support tool that combines geospatial analysis, graph neural networks, and user-centered interface design.